\documentclass[11pt,a4paper]{article}

\usepackage{geometry}
\usepackage{longtable}
\usepackage{array}
\usepackage{graphicx}
\usepackage{hyperref}

\geometry{margin=1in}

% ---------------- Customization ----------------
\newcommand{\AuthorName}{Elias Khalil}
\newcommand{\SemesterText}{Second Semester 2025/2026}
\newcommand{\LogoFile}{bzu_logo.png} % place logo in project folder
% ----------------------------------------------

\begin{document}

% ============================================================
% Front Page
% ============================================================
\begin{titlepage}
\begin{center}

\vspace*{1cm}

\includegraphics[width=0.35\textwidth]{\LogoFile}

\vspace{1.2cm}

\textbf{\Large Birzeit University}\\[0.4cm]
\textbf{\large Faculty of Engineering and Technology}\\[0.2cm]
\textbf{\large Department of Electrical and Computer Engineering}\\[0.6cm]

\textbf{\large \SemesterText}

\vfill

\textbf{\Large Course Project Specification}

\vspace{0.4cm}

\rule{3cm}{0.4pt}

\vspace{0.6cm}

\textbf{\Huge BIRD}\\[0.3cm]
{\Large \textit{Birzeit Integrated Router Design}}

\vfill

\textbf{\large Author: \AuthorName}

\vspace*{1cm}

\end{center}
\end{titlepage}

\newpage

% ============================================================
% Main Spec
% ============================================================
\title{\textbf{BIRD -- Birzeit Integrated Router Design}\\Functional Specification}
\author{}
\date{}
\maketitle

\tableofcontents
\newpage

% ============================================================
\section{Overview}
% ============================================================

BIRD (Birzeit Integrated Router Design) is a packet-based routing block that receives traffic from a single input interface and routes it to one of two output interfaces: \textbf{local} or \textbf{remote}. Routing decisions and packet processing behavior are determined by a 32-bit configuration signal (\texttt{cfg}) that accompanies packet data on the input interface.

BIRD supports two traffic types:
\begin{itemize}
  \item \textbf{Local traffic}: forwarded directly to the local output interface.
  \item \textbf{Remote traffic}: received as multiple fragments, accumulated, reordered, merged, and transmitted as a single packet.
\end{itemize}

BIRD does not signal errors. Any invalid or malformed packet is silently dropped and recorded internally.

% ============================================================
\section{Packet Data and Control Signals}
% ============================================================

\subsection{Packet Data Stream}

The packet data stream transferred on \texttt{data\_in} consists of:
\begin{verbatim}
Payload (1--255 bytes)
CRC16 (16 bits)
\end{verbatim}

\begin{itemize}
  \item Payload data is transferred at 8 bits per cycle.
  \item CRC16 covers only the payload bytes.
  \item CRC16 is transferred immediately after the payload.
\end{itemize}

\subsection{Configuration Signal (\texttt{cfg})}

The \texttt{cfg} signal is a 32-bit \textbf{sideband control signal} and is \textbf{not part of the packet data stream}.

\begin{itemize}
  \item \texttt{cfg} describes the fragment currently being transferred.
  \item \texttt{cfg} must be valid and stable during the transfer of each fragment.
  \item For remote traffic, \texttt{cfg} is provided with \textbf{each fragment}.
  \item For local traffic, \texttt{cfg} applies to the entire packet.
\end{itemize}

\subsection{Sampling Rule}

BIRD samples the \texttt{cfg} signal on the same clock cycle as the first payload byte of each fragment, when \texttt{in\_vld} and \texttt{in\_rdy} are asserted.

% ============================================================
\section{Interface Protocol}
% ============================================================

All interfaces use a valid/ready handshake.

\subsection{Transfer Rule}
A transfer occurs on a rising edge of \texttt{clk} when:
\[
\texttt{vld} = 1 \;\land\; \texttt{rdy} = 1
\]

\subsection{Stability Rule}
When \texttt{vld = 1} and \texttt{rdy = 0}, all associated data and control signals must remain stable.

% ============================================================
\section{Signal Interfaces}
% ============================================================

\subsection{Global Signals}

\begin{longtable}{|>{\ttfamily}p{3cm}|c|c|p{7cm}|}
\hline
Signal & Width & Direction & Description \\
\hline
clk & 1 & Input & System clock (rising-edge triggered) \\
rst\_n & 1 & Input & Active-low reset \\
\hline
\end{longtable}

\subsection{Status Output Interface}

\begin{longtable}{|>{\ttfamily}p{3cm}|c|c|p{7cm}|}
\hline
Signal & Width & Direction & Description \\
\hline
drop\_cnt & 16 & Output & Number of packets dropped (wrap-around counter) \\
\hline
\end{longtable}

\subsection{Input Interface (Producer $\rightarrow$ BIRD)}

\begin{longtable}{|>{\ttfamily}p{3cm}|c|c|p{7cm}|}
\hline
Signal & Width & Direction & Description \\
\hline
in\_vld & 1 & Input & Input data valid \\
in\_rdy & 1 & Output & BIRD ready to accept input \\
data\_in & 8 & Input & Payload and CRC data stream \\
cfg & 32 & Input & Sideband configuration for current fragment \\
\hline
\end{longtable}

\subsection{Local Output Interface}

\begin{longtable}{|>{\ttfamily}p{3cm}|c|c|p{7cm}|}
\hline
Signal & Width & Direction & Description \\
\hline
local\_vld & 1 & Output & Local output valid \\
local\_rdy & 1 & Input & Local consumer ready \\
data\_local & 8 & Output & Local payload data \\
\hline
\end{longtable}

\subsection{Remote Output Interface}

\begin{longtable}{|>{\ttfamily}p{3cm}|c|c|p{7cm}|}
\hline
Signal & Width & Direction & Description \\
\hline
remote\_vld & 1 & Output & Remote output valid \\
remote\_rdy & 1 & Input & Remote consumer ready \\
data\_remote & 32 & Output & Remote packet data stream \\
\hline
\end{longtable}

% ============================================================
\section{Configuration Word (\texttt{cfg}) Format}
% ============================================================

The \texttt{cfg} signal is a 32-bit sideband control word with the following format:

\begin{longtable}{|c|c|c|p{7cm}|}
\hline
Bits & Field & Width & Description \\
\hline
{[0]} & TRAFFIC\_TYPE & 1 & 0 = Local traffic, 1 = Remote traffic \\
{[7:1]} & Reserved & 7 & Must be 0 \\
{[15:8]} & PAYLOAD\_LEN & 8 & Payload length in bytes (1--255) \\
{[20:16]} & FRAG\_NUM & 5 & Fragment number within the packet (1--31) \\
{[23:21]} & Reserved & 3 & Must be 0 \\
{[28:24]} & SEQ\_NUM & 5 & Packet sequence number (1--31). 0 is invalid \\
{[31:29]} & Reserved & 3 & Must be 0 \\
\hline
\end{longtable}

% ============================================================
\section{Local Traffic Processing}
% ============================================================

A fragment is classified as local traffic when:
\[
\texttt{cfg[0]} = 0
\]

\begin{itemize}
  \item Local traffic consists of a single fragment only.
  \item Payload bytes are forwarded directly to the local output interface.
  \item No buffering, accumulation, or reordering is performed.
  \item \texttt{FRAG\_NUM} shall be equal to 1.
  \item \texttt{SEQ\_NUM} identifies the packet but has no functional impact on local routing.
  \item The input CRC16 is forwarded unchanged.
\end{itemize}

% ============================================================
\section{Remote Traffic Processing}
% ============================================================

A fragment is classified as remote traffic when:
\[
\texttt{cfg[0]} = 1
\]

\subsection{Fragment Rules}

\begin{itemize}
  \item A remote packet consists of one or more fragments.
  \item Each fragment carries a \texttt{FRAG\_NUM} field indicating its position within the packet.
  \item All fragments belonging to the same packet share the same \texttt{SEQ\_NUM}.
  \item Valid fragment numbers are $1,2,\dots,N$, where $N$ is determined by received fragments.
  \item Fragments may arrive out of order.
  \item Only one remote packet (identified by \texttt{SEQ\_NUM}) may be accumulated at a time.
\end{itemize}

\subsection{Accumulation and Reassembly}

\begin{enumerate}
  \item Each fragment payload is buffered and indexed by \texttt{FRAG\_NUM}.
  \item All fragments sharing the same \texttt{SEQ\_NUM} are considered part of the same packet.
  \item When all required fragments are received, fragments are reordered based on \texttt{FRAG\_NUM}.
  \item Payloads are concatenated to form a merged payload.
  \item A new CRC16 is generated over the merged payload.
\end{enumerate}

\subsection{Remote Output Format}

The remote output interface transmits:
\begin{verbatim}
Merged Payload
CRC16
\end{verbatim}

The \texttt{cfg} signal is not transmitted on the remote output interface.

% ============================================================
\section{Silent Packet Drop Policy}
% ============================================================

BIRD does not signal errors. Invalid packets are silently dropped and counted.

\subsection{Drop Conditions}

A packet is dropped if any of the following occur:
\begin{itemize}
  \item \texttt{SEQ\_NUM == 0}
  \item \texttt{FRAG\_NUM == 0}
  \item \texttt{PAYLOAD\_LEN} is outside the range 1--255
  \item Any reserved bits in \texttt{cfg} are non-zero
  \item A fragment arrives with a mismatched \texttt{SEQ\_NUM} while another packet is being accumulated
  \item A required fragment for a packet is missing
  \item A fragment with \texttt{FRAG\_NUM == 1} arrives while a previous packet with a different \texttt{SEQ\_NUM} is still incomplete
\end{itemize}

\subsection{Drop Counter Behavior}

\begin{itemize}
  \item \texttt{drop\_cnt} increments by one for each packet that is dropped.
  \item The counter increments once per packet, regardless of the number of fragments received.
  \item \texttt{drop\_cnt} is a modulo-$2^{16}$ counter and wraps around on overflow.
  \item \texttt{drop\_cnt} is not affected by successful packet transmission.
\end{itemize}

When a drop condition is detected, all buffered data for the affected packet is immediately discarded.

% ============================================================
\section{Reset Behavior}
% ============================================================

When \texttt{rst\_n = 0}:
\begin{itemize}
  \item All valid outputs are deasserted.
  \item All internal buffers and fragment state are cleared.
  \item Any packet in progress is discarded.
  \item \texttt{drop\_cnt} is cleared to zero.
\end{itemize}

Normal operation resumes after reset deassertion.

% ============================================================
\section{Notes for Verification}
% ============================================================

\begin{itemize}
  \item Verify correct classification of local versus remote traffic.
  \item Verify handshake stability under backpressure.
  \item Verify correct use of \texttt{FRAG\_NUM} and \texttt{SEQ\_NUM}.
  \item Verify fragment reordering and merged payload correctness.
  \item Verify silent drop behavior under all listed conditions.
  \item Verify correct incrementation and wrap-around behavior of \texttt{drop\_cnt}.
  \item Verify CRC16 regeneration for remote packets.
\end{itemize}

\end{document}
